\chapter{函数空间 \texorpdfstring{$C(K)$}{C(K)}}
\label{chap:continuous-function-space}

此前讨论函数列时，函数只是带有自变量的表达式；一致收敛定理却已经暗示，
一整个连续函数也可以作为空间中的一个点。两个函数之间的距离取为最大偏差，
逐点变化的问题便转化为这个空间中的收敛、完备性和紧致性问题。

本章设 $K$ 为非空紧致 Hausdorff 空间，标量域为 $\mathbb F=\R$ 或 $\C$。
若只关心第一遍阅读，可以把 $K$ 始终看成闭区间或紧致度量空间。

\section{上确界范数与一致收敛}

记 $C(K;\mathbb F)$ 为 $K$ 上所有连续 $\mathbb F$ 值函数组成的线性空间。连续函数在
紧集上有界并取得模的最大值，因而
\[
 \norm f_\infty=\max_{x\in K}|f(x)|
\]
总是有限的。

\begin{proposition}
\label{prop:sup-norm-on-ck}
$\norm{\cdot}_\infty$ 是 $C(K;\mathbb F)$ 上的范数，并且
\[
 \norm{fg}_\infty\le\norm f_\infty\norm g_\infty.
\]
由该范数产生的收敛恰为 $K$ 上的一致收敛。
\end{proposition}

\begin{proof}
正定性、齐次性和三角不等式逐点成立，取最大值后仍成立；乘积不等式同理。
而 $\norm{f_n-f}_\infty<\varepsilon$ 等价于对所有 $x\in K$ 同时有
$|f_n(x)-f(x)|<\varepsilon$，这正是一致收敛定义中的量词次序。
\end{proof}

完备性是这个函数空间的第一个重要性质。它把
\chapterref{chap:uniform-convergence-preservation}“连续函数的一致极限连续”
改写成空间本身没有缺口。

\begin{theorem}[\texorpdfstring{$C(K)$}{C(K)} 的完备性]
\label{thm:ck-banach}
赋予上确界范数后，$C(K;\mathbb F)$ 是 Banach 空间。
\end{theorem}

\begin{proof}
设 $(f_n)$ 是上确界范数下的 Cauchy 列。对每个 $x\in K$，
$|f_n(x)-f_m(x)|\le\norm{f_n-f_m}_\infty$，所以 $(f_n(x))$ 是 $\mathbb F$ 中的
Cauchy 列。标量域完备，故可定义
$f(x)=\lim_n f_n(x)$。

给定 $\varepsilon>0$，取 $N$ 使 $m,n\ge N$ 时
$\norm{f_n-f_m}_\infty<\varepsilon$。固定 $n\ge N$ 并令 $m\to\infty$，
逐点得到 $|f_n(x)-f(x)|\le\varepsilon$，且此式对所有 $x$ 成立。因此
$\norm{f_n-f}_\infty\le\varepsilon$，即 $f_n$ 一致收敛于 $f$。连续函数的
一致极限连续，故 $f\in C(K;\mathbb F)$。所以每个 Cauchy 列都在该空间中收敛。
\end{proof}

证明中有两个完备性层次：先用 $\R$ 或 $\C$ 的完备性逐点定义候选极限，
再用 Cauchy 估计的一致性把逐点极限提升为一致极限。若把上确界范数换成
只测量平均误差的量，第二步一般不再成立。

\section{积分范数与不同的接近方式}

在 $C([a,b])$ 上还可以定义
\[
 \norm f_p=\left(\int_a^b|f(x)|^p\,dx\right)^{1/p},\qquad 1\le p<\infty.
\]
它们也是范数。这里把三角不等式所需的有限维事实写出，以免提前调用第三卷
的 \(L^p\) 理论。设 \(1<p<\infty\) 且 \(1/p+1/q=1\)。Young 不等式
\[
 uv\le \frac{u^p}{p}+\frac{v^q}{q}\qquad(u,v\ge0)
\]
经齐次化给出有限和的 Hölder 不等式
\[
 \sum_i|a_ib_i|
 \le\left(\sum_i|a_i|^p\right)^{1/p}
    \left(\sum_i|b_i|^q\right)^{1/q}.
\]
把它分别用于
\(|a_i|\,|a_i+b_i|^{p-1}\) 和
\(|b_i|\,|a_i+b_i|^{p-1}\)，得到
\[
 \left(\sum_i|a_i+b_i|^p\right)^{1/p}
 \le
 \left(\sum_i|a_i|^p\right)^{1/p}
 +\left(\sum_i|b_i|^p\right)^{1/p};
\]
若左端为零则无需相除，非零时约去
\((\sum_i|a_i+b_i|^p)^{1/q}\) 即可。\(p=1\) 时直接使用数的三角
不等式。对同一分割上的标记值，把 \(a_i,b_i\) 分别换成
\(f(\xi_i)(\Delta x_i)^{1/p}\)、
\(g(\xi_i)(\Delta x_i)^{1/p}\)，再令 Riemann 和趋于积分，便得到
\[
 \norm{f+g}_p\le\norm f_p+\norm g_p.
\]
这个论证适用，是因为连续函数的有关幂次都 Riemann 可积。第三卷将在
可测函数上重新证明一般的 Hölder 与 Minkowski 不等式。
由 $|f(x)|\le\norm f_\infty$ 立即得到
\begin{equation}
 \norm f_p\le(b-a)^{1/p}\norm f_\infty.
 \label{eq:sup-controls-integral-norm}
\end{equation}
所以一致收敛必推出每个积分范数下的收敛。

反向推论不成立。令 $f_n$ 是以 $1/2$ 为中心、高为 $1$、底宽为 $2/n$ 的
三角形函数。则 $\norm{f_n}_\infty=1$，而
$\norm{f_n}_p^p=2/(n(p+1))\to0$。积分范数允许误差集中在越来越小的区间；
上确界范数只看最大误差，不考虑其发生区域有多小。

更深的差别出现在完备性上。

\begin{proposition}
\label{prop:c-interval-not-complete-l1}
$C([0,1])$ 赋予 $\norm{\cdot}_1$ 后不是完备空间。
\end{proposition}

\begin{proof}
令 $f_n$ 在 $[0,1/2-1/n]$ 上为 $0$，在 $[1/2+1/n,1]$ 上为 $1$，中间作
线性连接。它在 $\norm{\cdot}_1$ 下趋于阶跃函数
$h=\mathbf1_{(1/2,1]}$，因二者只在长度 $2/n$ 的区间上可能不同。故
$(f_n)$ 是 $\norm{\cdot}_1$-Cauchy 列。

若它在 $C([0,1])$ 中还收敛于某个连续函数 $g$，则三角不等式给出
$\norm{g-h}_1=0$。连续性于是迫使 $g=0$ 于 $[0,1/2)$、$g=1$ 于
$(1/2,1]$，与 $g$ 在 $1/2$ 连续矛盾。因此该 Cauchy 列在连续函数空间中
没有 $\norm{\cdot}_1$ 极限。
\end{proof}

第三卷将通过把几乎处处相等的可测函数视为同一元素，补入这类积分范数下
缺失的极限，得到完备的 $L^p$ 空间。

\section{点值泛函与积分泛函}

为了比较点值与积分这两类映射，只需用到有界线性算子的定义和最基本的
连续性判据。这里不展开算子空间、对偶空间的完备性或算子之间的运算；这些
内容将在泛函分析中系统建立。

若 $X$、$Y$ 是赋范空间，线性映射 $T:X\to Y$ 在零点连续，当且仅当存在
$C\ge0$ 使 $\norm{Tx}\le C\norm x$。这样的线性映射称为有界线性算子；
最小的可用常数是算子范数
\[
 \norm T=\sup_{\norm x\le1}\norm{Tx}.
\]

线性性使“在零点连续”已经包含所有点的连续性。若
$\norm{Tx}\le C\norm x$，则
\[
 \norm{Tx-Tx_0}\le C\norm{x-x_0}.
\]
反之，若 $T$ 在零点连续，可取 $r>0$，使 $\norm x<r$ 时
$\norm{Tx}<1$。对 $x\ne0$，把 $x$ 缩放为 $rx/(2\norm x)$，便得
\[
 \norm{Tx}\le\frac2r\norm x.
\]
因此连续性不只是定性条件，它给出一个全局线性增长界；算子范数就是所有
这类增长常数的下确界。

对 $x_0\in K$，点值映射
$\delta_{x_0}(f)=f(x_0)$ 是 $C(K)$ 上的线性泛函，而且
\[
 |\delta_{x_0}(f)|\le\norm f_\infty,
 \qquad \norm{\delta_{x_0}}=1.
\]
等号由常函数 $1$ 取得。因而上确界范数收敛时，每一点的函数值都连续地
随函数变化。

若 $K=[a,b]$，积分映射
\[
 I(f)=\int_a^bf(x)\,dx
\]
也是有界线性泛函，且 $\norm I=b-a$。更一般地，对
$g\in C([a,b])$，
\[
 I_g(f)=\int_a^bf(x)g(x)\,dx,
 \qquad \norm{I_g}=\int_a^b|g(x)|\,dx.
\]
上界来自积分不等式。为得到反向不等式，可用连续函数
$f_\eta=\overline g/\sqrt{|g|^2+\eta^2}$ 逼近 $g$ 的共轭相位；在实值情形
这就是符号函数的连续近似。令 $\eta\downarrow0$，
$\norm{f_\eta}_\infty\le1$，并且对 $r\ge0$ 有
\[
 0\le r-\frac{r^2}{\sqrt{r^2+\eta^2}}\le\eta.
\]
所以
\[
 0\le\int_a^b|g|-I_g(f_\eta)\le(b-a)\eta.
\]
令 $\eta\downarrow0$ 即得 $\norm{I_g}\ge\int|g|$，从而范数公式成立。

点值泛函在积分范数下却一般不连续。前面的三角形函数满足
$f_n(1/2)=1$ 而 $\norm{f_n}_1\to0$。同一个代数映射是否连续，取决于空间
采用的范数。

\section{函数代数与格运算}

$C(K;\mathbb F)$ 不仅是线性空间，还是含单位元的交换代数：逐点乘积仍连续，且
命题 \ref{prop:sup-norm-on-ck} 给出范数的次乘性。若标量为实数，还可逐点
比较函数，定义
\[
 f\vee g=\max\{f,g\},\qquad f\wedge g=\min\{f,g\}.
\]
恒等式
\[
 f\vee g=\frac{f+g+|f-g|}{2},\qquad
 f\wedge g=\frac{f+g-|f-g|}{2}
\]
说明二者仍属于 $C(K;\R)$。并且
\[
 \norm{|f|-|g|}_\infty\le\norm{f-g}_\infty,
\]
所以取绝对值、最大值和最小值在上确界范数下都是连续运算。

这些格运算在\chapterref{chap:polynomial-approximation}的
Stone--Weierstrass 定理中承担关键角色：分离两个点
容易构造一个局部合适的函数，有限次最大、最小运算则把局部信息拼成整个
紧集上的一致逼近。

\section{闭有界集为何不再紧致}

有限维空间中，闭且有界等价于紧致。函数空间通常是无限维的，这个结论随即
失效。下面的引理给出普遍原因。

\begin{lemma}[Riesz 引理]
\label{lem:riesz-lemma}
设 $M$ 是赋范空间 $X$ 的真闭线性子空间。对每个 $0<\alpha<1$，存在
$x\in X$，使 $\norm x=1$ 且
\[
 \inf_{y\in M}\norm{x-y}>\alpha.
\]
\end{lemma}

\begin{proof}
取 $z\notin M$。因 $M$ 闭，$d=\inf_{y\in M}\norm{z-y}>0$。选
$y_0\in M$ 使 $\norm{z-y_0}<d/\alpha$，并令
$x=(z-y_0)/\norm{z-y_0}$。任取 $y\in M$，有
\[
 \norm{x-y}=\frac{\norm{z-(y_0+\norm{z-y_0}y)}}{\norm{z-y_0}}
 \ge\frac d{\norm{z-y_0}}>\alpha,
\]
因为括号中的点属于 $M$。
\end{proof}

\begin{theorem}
\label{thm:unit-ball-compact-iff-finite-dimensional}
赋范空间的闭单位球紧致，当且仅当该空间有限维。
\end{theorem}

\begin{proof}
有限维情形由 Heine--Borel 定理成立。反之，若 $X$ 无限维，递归令
$M_n=\operatorname{span}\{x_1,\ldots,x_n\}$，并用 Riesz 引理选单位向量
$x_{n+1}$，使它到 $M_n$ 的距离大于 $1/2$。于是 $m<n$ 时
$\norm{x_n-x_m}>1/2$，所以闭单位球中的序列 $(x_n)$ 没有 Cauchy 子列，
更没有收敛子列；闭单位球不紧致。
\end{proof}

例如 $1,x,x^2,\ldots$ 在 $C([0,1])$ 中线性无关，因此该空间无限维；它的
闭单位球虽闭且有界，却不紧。要从一个函数族中抽出一致收敛子列，必须补充
约束函数随自变量变化的条件。下一章的等度连续性正是所需条件。

也可以在 $C([0,1])$ 中直接看见这种失效。取两两不交、端点趋于 $0$ 的
闭区间 $I_n\subset(0,1]$，在每个 $I_n$ 上作高为 $1$、端点为 $0$ 的
三角形帽函数 $h_n$，并在区间外令它为零。则
$\norm{h_n}_\infty=1$，而 $m\ne n$ 时
$\norm{h_n-h_m}_\infty=1$。因此 $(h_n)$ 没有 Cauchy 子列。所有函数都
具有同一个值域界，但峰的位置可以不断移动并缩窄；下一章将用等度连续性
排除这种局部变化。

\section*{习题}
\addcontentsline{toc}{section}{习题}

\subsection*{A组　范数与完备性}
\begin{enumerate}[label=\textbf{\thechapter.\arabic*.}]
 \exerciseitem{ck-sup-norm-details}
 完整验证上确界范数的范数公理，并说明紧致性在定义中何处使用。
 \exerciseitem{ck-uniform-cauchy-criterion}
 证明函数列在 $C(K)$ 中收敛，当且仅当它在上确界范数下为 Cauchy 列。
 \exerciseitem{integral-norm-triangle-minkowski-finite}
 不引用后文 $L^p$ 理论，证明 $\norm{\cdot}_1$ 和 $\norm{\cdot}_2$ 的三角不等式。
 \exerciseitem{spike-integral-norms}
 计算正文三角形函数的 $p$-积分范数，并说明它不逐点收敛于零的点。
 \exerciseitem{c-not-complete-lp}
 把命题 \ref{prop:c-interval-not-complete-l1} 推广到每个 $1\le p<\infty$。
\end{enumerate}

\subsection*{B组　泛函与代数结构}
\begin{enumerate}[label=\textbf{\thechapter.\arabic*.},resume]
 \exerciseitem{evaluation-functional-norm}
 求有限个点值线性组合 $T(f)=\sum_{j=1}^m a_jf(x_j)$ 的范数，其中 $x_j$
 两两不同。
 \exerciseitem{weighted-integral-functional-norm}
 补全 $\norm{I_g}=\int|g|$ 的证明，并解释当 $g$ 有零点时为何不能直接取
 $f=\operatorname{sgn}g$。
 \exerciseitem{point-evaluation-l1-discontinuous}
 构造连续函数列，证明点值泛函在 $\norm{\cdot}_1$ 下不连续。
 \exerciseitem{ck-lattice-lipschitz}
 证明 $f\mapsto|f|$ 是 $1$-Lipschitz，并估计最大值、最小值运算的连续性。
 \exerciseitem{neumann-inverse-function-algebra}
 若 $\norm f_\infty<1$，证明 $1-f$ 在 $C(K)$ 中可逆，并证明
 $(1-f)^{-1}=\sum_{n=0}^\infty f^n$ 在上确界范数下收敛。
\end{enumerate}

\subsection*{C组　维数与紧致性}
\begin{enumerate}[label=\textbf{\thechapter.\arabic*.},resume]
 \exerciseitem{ck-infinite-dimensional}
 对无限紧致 Hausdorff 空间 $K$，证明 $C(K)$ 无限维。
 \exerciseitem{disjoint-tent-no-uniform-subsequence}
 在 $C([0,1])$ 的闭单位球中显式构造两两距离为 $1$ 的函数序列。
 \exerciseitem{finite-dimensional-unit-ball-compact}
 详细证明有限维赋范空间的闭单位球紧致，不预先假定范数等价。
 \exerciseitem{riesz-lemma-closedness}
 给出非闭子空间的例子，说明 Riesz 引理中的闭性不能删除。
 \optionalexerciseitem{compact-operator-first-example}
 设 $k\in C([0,1]^2)$，$Tf(x)=\int_0^1k(x,y)f(y)\,dy$。证明 $T$ 有界，
 并先猜测为何它会把单位球送入比一般有界集更紧的函数族。
\end{enumerate}
